% ==================================================================
% id:       cond-016
% topic:    הסתברות מותנית ואי-תלות
% subtopic: נוסחת ההסתברות השלמה ונוסחת בייס
% solution: full
% ==================================================================
במבחן מסוים מצליחים $\frac23$ מהנבחנים. רק מחצית מהנבחנים יודעים את החומר, ו-$\frac34$ מהיודעים מצליחים. מה ההסתברות שנבחן שלא יודע את החומר יצליח?
\hint{כתבו את נוסחת ההסתברות השלמה עבור ״מצליח״. הנעלם הוא ההסתברות שנבחן מצליח, בהינתן שהוא לא יודע את החומר.}
\begin{solbox}
נסמן $K$ = ״הנבחן יודע את החומר״, $S$ = ״הנבחן מצליח״. הנתונים:
\[
\PP(S)=\tfrac23,\qquad \PP(K)=\tfrac12,\qquad \PP(S\mid K)=\tfrac34,
\]
ומחפשים את $x=\PP(S\mid K^C)$.

המאורעות $K$ ו-$K^C$ הם חלוקה של $\Omega$, ולכן לפי נוסחת ההסתברות השלמה
\[
\PP(S)=\PP(S\mid K)\,\PP(K)+\PP(S\mid K^C)\,\PP(K^C).
\]
נציב את הנתונים, ו-$\PP(K^C)=1-\frac12=\frac12$:
\[
\tfrac23=\tfrac34\cdot\tfrac12+x\cdot\tfrac12=\tfrac38+\tfrac x2.
\]
מכאן $\frac x2=\frac23-\frac38=\frac7{24}$, ולכן $x=\ans{\frac7{12}}$.
\end{solbox}
